% Copyright 2026 The terCAD team. All rights reserved.
% Use of this content is governed by a CC BY-NC-ND 4.0 license that can be found in the LICENSE file.

\subsection{Looking Back}
\markboth{Continuing}{Looking Back}

Throughout this book, we have traced the journey of money in Fingrom, from allocating it to
budget categories and maintaining accounts, to recording every bill, income and transfer, and
directing savings towards goals. Together, these steps form a complete cycle of personal
financial management.

\begin{enumerate}
  \item \textbf{Plan} -- define budget categories and their limits, set goals with deadlines;
  \item \textbf{Record} -- enter transactions as they happen, automate the recurring ones,
  import the history;
  \item \textbf{Review} -- check the Summary table, the metrics, and the progress of goals;
  \item \textbf{Adjust} -- correct the limits, the deadlines, and the habits, and start the
  next cycle.
\end{enumerate}

\noindent No application can make financial decisions on behalf of a person. However, it can
make the consequences of these decisions visible quickly, honestly and effortlessly. This is
precisely why Fingrom was created: to be an intuitive, efficient and inclusive tool that
enables anyone to manage their finances and leaves no one behind.


\subsection{Work in Progress}
\markboth{Continuing}{Work in Progress}

Fingrom is in an active state of development. Since the first code was committed in 2023, the
application has grown from a prototype developed in a developer's spare time into a product
published in all major stores, with new releases appearing regularly. Consequently:

\begin{itemize}
  \item the appearance of some pages may differ from the descriptions and screenshots of this
  book, especially in other languages or on other screen sizes;
  \item new features are added, and existing ones are refined, based on the feedback of users;
  \item some behaviors described here may change, if a more convenient or more correct
  approach is found (but that book will be updated accordingly).
\end{itemize}

\noindent The source code is the ultimate reference, and it is open to everyone at
\href{https://github.com/lyskouski/app-finance}{https://github.com/lyskouski/app-finance}. The
release notes describe what has changed in each version. Even that book is a part of that
repository (\q{./docs/user-manual/README.md}).


\subsection{Comparing with Other Solutions}
\markboth{Continuing}{Comparing with Other Solutions}

The financial applications market is large, and there is no single best choice for everyone.
The most common alternatives can be divided into a few categories:

\begin{description}
  \item[Desktop double-entry software] (e.g. GnuCash, KMyMoney) -- free and open-source
  programs that follow the full accounting model with assets, liabilities, equity, incomes,
  and expenses. They are powerful and precise, suitable even for small businesses, but demand
  knowledge of accounting and are designed primarily for desktop computers.
  \item[Local-first budgeting tools] (e.g. HomeBank, Actual Budget) -- focus on budgeting and
  keep data under the control of the user. Some of them require a self-hosted server for
  synchronization between devices.
  \item[Self-hosted web applications] (e.g. Firefly III) -- rich functionality accessible
  from any browser, at the cost of installing and maintaining one's own server.
  \item[Commercial cloud services] (e.g. YNAB, Quicken) -- polished applications with direct
  connections to banks and automatic import of transactions, usually by subscription, and
  with the financial history stored on the servers of the provider. The shutdown of Mint in
  2024 showed the other side of this convenience: when a service closes, its users have to
  migrate years of history elsewhere.
\end{description}

\noindent Fingrom occupies a unique position between these groups: it is as private as
local-first tools, as accessible as cloud services and far simpler than double-entry software
-- all while remaining free.

\begin{center}
\footnotesize
\begin{tabularx}{\textwidth}{|l|X|X|X|}
  \hline
  \textbf{Criterion} & \textbf{Fingrom} & \textbf{Desktop double-entry} & \textbf{Commercial cloud} \\
  \hline
  Cost & free & free & subscription \\
  Ads & none & none & none or some \\
  Platforms & all platforms & mostly desktop & all platforms \\
  Data storage & on own devices & on own device & on servers \\
  Serverless sync & yes (P2P) & no & -- \\
  Bank connection & no & limited & yes \\
  Double-entry & no & yes & no \\
  Learning curve & low & high & low \\
  \hline
\end{tabularx}
\end{center}


\subsection{Strengths and Limitations}
\markboth{Continuing}{Strengths and Limitations}

A fair manual should not only demonstrate the strengths of an application, but also highlight
its limitations. This information is useful for deciding whether Fingrom meets a particular
need and where to look for complementary functionality.

\subsubsection{Strengths}

\begin{itemize}
  \item \textbf{Free, without ads, open source.} No feature is locked behind a subscription,
  and anyone can verify how the application works.
  \item \textbf{Privacy by design.} No registration is required; the financial data is stored
  locally, can be encrypted, and travels only between the devices of the user. Anonymous
  usage statistics (device type, language, visited pages, errors) are collected solely to
  improve the application, as described in the privacy policy.
  \item \textbf{Every platform.} iOS, Android, macOS, Windows, Linux, and the Web, with the
  same features everywhere.
  \item \textbf{Inclusive interface.} More than twenty languages, including right-to-left
  ones, adaptive layouts from phones to wide monitors, keyboard navigation, and zoom.
  \item \textbf{Multi-currency.} Accounts, budgets, and goals in any currency or
  cryptocurrency, with conversion of totals into the default one.
  \item \textbf{Flexible budgeting.} Fixed, relative, and unlimited limits, monthly
  coefficients, and weekly, monthly, or yearly intervals.
  \item \textbf{Insights.} A forecast of the budget based on Monte Carlo simulation, the
  Summary table, charts, and the automatic progress of goals with a monthly target.
  \item \textbf{Low effort.} Category suggestions, remembered choices, recurring payments,
  and import from \q{CSV}, \q{QIF}, and \q{OFX}.
\end{itemize}

\subsubsection{Limitations}

\begin{itemize}
  \item \textbf{No direct bank connection.} Transactions are entered manually, created by
  recurring payments, or imported from files exported by the bank.
  \item \textbf{Single-entry model.} Fingrom is designed for personal and family finances,
  not for formal bookkeeping: it does not produce statutory reports, does not track
  receivables and payables, and does not replace an accountant for business or tax purposes.
  \item \textbf{No investment portfolio.} Securities can be tracked only as the balance of a
  deposit account, without quotes, positions, or returns.
  \item \textbf{No roles and permissions.} All synchronized devices have equal access to all
  data; there is no read-only mode for a family member or an advisor.
  \item \textbf{Limited import.} Transfers are not imported from files, and the mapping of
  columns has to be repeated for every file.
  \item \textbf{Manual reconciliation.} There is no mark for transactions verified against a
  bank statement; discrepancies are corrected via the balance of an account.
\end{itemize}


\subsection{Possible Improvements}
\markboth{Continuing}{Possible Improvements}

From the outset, the concept of Fingrom encompassed more ideas than could be implemented
immediately. Some of these features are already included in the application, such as adaptive
design, onboarding, personalization, encryption, multi-factor protection, charts and
automatic categorization. Others remain on the horizon.

\begin{description}
  \item[Integration with banks] -- connecting accounts through Open Banking interfaces, with
  automatic import of transactions and real-time balances, while keeping the data under the
  control of the user.
  \item[Reconciliation] -- marking transactions as verified against a bank statement, to
  identify discrepancies and missing entries.
  \item[Intelligent notifications] -- detecting anomalies such as unexpectedly high expenses,
  possible duplicate or unauthorized charges, and suggesting adjustments of budget limits
  based on actual spending.
  \item[Receipt recognition] -- extracting the amount, the date, and the merchant from a
  photo of a receipt via optical character recognition (OCR).
  \item[Parsing notifications] -- creating transactions from the messages of banking
  applications on the device.
  \item[Reports] -- customizable report templates and periodic summaries, exported to common
  formats.
  \item[Shared access] -- roles for family members or advisors with restricted permissions.
  \item[Investments] -- tracking securities with their quotes and returns.
\end{description}

\noindent The order in which these ideas are implemented depends largely on the users.
Development priorities are driven by feedback received through the issue tracker, stores and
direct messages.


\subsection{Getting Involved}
\markboth{Continuing}{Getting Involved}

Fingrom is an open project, and there are many ways to influence its future, even without
programming:

\begin{itemize}
  \item \textbf{Report issues and ideas} at
  \href{https://github.com/lyskouski/app-finance/issues}{https://github.com/lyskouski/app-finance/issues}
  -- each report helps to find a problem or to prioritize a feature;
  \item \textbf{Translate} the interface, or improve an existing translation;
  \item \textbf{Test} pre-release versions via TestFlight, Google Play, or the web version;
  \item \textbf{Contribute} code or documentation, including this book, via pull requests;
  \item \textbf{Rate and recommend} the application in the stores -- it helps others to find
  it;
  \item \textbf{Sponsor} the development: subscribing does not unlock any features, but
  supports the evolution of the application and preserves its availability.
\end{itemize}

\noindent Every contribution is acknowledged in the \textbf{About} -- \textbf{Contributors}
section of the application.

\vspace{1cm}

\noindent Thank you for reading this book. May your budgets balance, your accounts grow, and
your goals come true.
